% !PRINTED 70
\weberpage{70}
%<<CONT

Nennen wir $\Theta$-Funktionen mit denselben Perioden gleicher
Ordnung und gleicher Charakteristik \gesperrt{verwandt}, und bezeichnen
wir ferner ein System von Funktionen als \gesperrt{linear abh\"angig}
oder \gesperrt{unabh\"angig}, je nachdem eine homogene lineare Relation
mit konstanten (nicht s\"amtlich verschwindenden) Koeffizienten
zwischen diesen Funktionen besteht oder nicht besteht, so gilt
der Satz:

\gesperrt{$m + 1$ verwandte $\Theta$-Funktionen $m$ter Ordnung sind
immer linear abh\"angig;}

oder:

\gesperrt{Aus $m$ linear unabh\"angigen verwandten $\Theta$-Funktionen
 $m$ter Ordnung l\"a\ss t sich jede andere verwandte
$\Theta$-Funktion linear (mit konstanten Koeffizienten)
zusammensetzen.}

Dieser Satz ist f\"ur unsere folgenden Betrachtungen von der
fundamentalsten Bedeutung; es ergibt sich daraus nach \S\,18, 3.,
da\ss\ jede $T$-Funktion $m$ter Ordnung mit Anwendung der Formel (3)
aus $m$ linear unabh\"angigen $\Theta$-Funktionen zusammengesetzt
werden kann.

Nach \S\,17, (11) ist die Summe der $m$ Argumentwerte, f\"ur
welche die Funktion $\Theta_{g_1,\, g_2}(v)$ im Periodenparallelogramm verschwindet
, nach dem Modul $(1, \omega)$ kongruent mit
\[
m\,\frac{1+\omega}{2} + \frac{g_1\omega - g_2}{2}.
\]

Die bis jetzt gegebene Definition der Funktion $\vartheta$ l\"a\ss t einen
Faktor, der eine Funktion von $\omega$ sein kann, unbestimmt. Wir
k\"onnen aber durch einen Zusatz zur Definition diesen Faktor
noch n\"aher bestimmen.

Wenn man die Definitionsgleichungen (1) zweimal nach $v$
und einmal nach $\omega$ differentiiert, indem man $g_1, g_2$ als konstant
ansieht, so ergibt sich nach einfacher Rechnung, da\ss\ die Funktion
\[
\frac{\partial^2 \vartheta(v,\omega)}{\partial v^2} - 4\pi i\,\frac{\partial \vartheta(v,\omega)}{\partial \omega}
\]
selbst den Bedingungen (1) gen\"ugt und also die Form
\[
\varphi(\omega)\,\vartheta(v,\omega)
\]
hat, worin $\varphi(\omega)$ in bezug auf $v$ konstant, also eine Funktion von $\omega$
allein ist. Wenn man jetzt
\[
e^{-\frac{1}{4\pi i}\int \varphi(\omega)\,d\omega}\ \vartheta(v,\omega)
\]
%CONT>>
